% Generated draft: requires mathematical and source review.
\documentclass[11pt]{article}
\usepackage[margin=25mm]{geometry}
\usepackage{fontspec}
\setmainfont{Latin Modern Roman}
\usepackage{amsmath,amssymb,mathtools,unicode-math}
\setmathfont{Latin Modern Math}
\usepackage{enumitem,needspace,xurl,hyperref}
\setlength{\parindent}{0pt}
\setlength{\parskip}{5pt}
\setlength{\emergencystretch}{3em}
\newcommand{\N}{\mathbb{N}}
\newcommand{\Z}{\mathbb{Z}}
\newcommand{\Q}{\mathbb{Q}}
\newcommand{\R}{\mathbb{R}}
\newcommand{\C}{\mathbb{C}}
\newcommand{\F}{\mathbb{F}}
\newcommand{\A}{\mathbb{A}}
\newcommand{\PP}{\mathbb{P}}
\newcommand{\E}{\mathbb{E}}
\newcommand{\eps}{\varepsilon}
\newcommand{\dd}{\,\mathrm{d}}
\begin{document}
\section*{UnsolvedMath ID: 1101216}
\textbf{Source problem code:} catalogue entry
\subsection{Definitions and mathematical statement}
Let $F_n$ be the free group of finite rank $n$, and let $h\in\operatorname{Aut}(F_n)$. For a finite-index subgroup $K\leq F_n$ satisfying $h(K)=K$, write $K'=[K,K]$ for its commutator subgroup and
\[
 h_{K,\mathrm{ab}}:K/K'\longrightarrow K/K',\qquad xK'\longmapsto h(x)K'.
\]
Eigenvalues mean eigenvalues over $\C$ of the induced linear automorphism of $(K/K')\otimes_{\Z}\C$. A complex number $\zeta$ is a root of unity if $\zeta^m=1$ for some integer $m\geq1$.

\textbf{Question.} For every automorphism $h$ of $F_n$, does there exist $K\leq F_n$ such that
\[
 [F_n:K]<\infty,\qquad h(K)=K,\qquad
 \exists m\geq1\ \exists v\in(K/K')\otimes_{\Z}\C\setminus\{0\}:
 (h_{K,\mathrm{ab}}\otimes1)^m v=v?
\]
Thus the requested eigenvalue may become $1$ after passing to a positive power of the induced automorphism. The source does not state the rank range. The usual positive-rank interpretation is $n\geq1$; a universal statement allowing $n=0$ has no nonzero abelianization vector and needs separate clarification.

For the supplied progress discussion, define the mapping-torus group and virtual first Betti number by
\[
 G_h=F_n\rtimes_h\Z,\qquad
 \operatorname{vb}_1(G_h)=\sup_{[G_h:L]<\infty}\dim_{\Q}H_1(L;\Q).
\]
The registry labels the question solved and claims a positive answer for every $h$ and every $F_n$, using an asserted equivalence with $\operatorname{vb}_1(G_h)\geq2$. It further claims $\operatorname{vb}_1(G_h)=\infty$ for $n\geq2$, citing Brinkmann's hyperbolicity criterion, Button's largeness theorem in the $\Z^2$ case, and Hagen--Wise together with Agol in the hyperbolic case. This is retained as a supplied synthesis: its author explicitly reports finding no publication drawing this precise corollary, and says that the combination was checked by that author. This draft does not independently verify the equivalence, the case split, or that combination.
\subsection{Short English statement}
Does every free-group automorphism preserve a finite-index subgroup on whose abelianization it has a root-of-unity eigenvalue? The registry reports a positive answer assembled from published ingredients, while expressly stating that no publication of this precise corollary was found.
\subsection{Sources}
\begin{itemize}
\item Local UnsolvedMath registry, permanent ID 1101216, source code catalogue entry, ``Questions in Geometric Group Theory --- Q 12.16''; status solved, classification SOLVED-IN-LITERATURE, updated 2026-08-06.
\item Mladen Bestvina, \emph{Questions in Geometric Group Theory}, list dated 2004, Question 12.16, PDF page 22, \url{https://www.math.utah.edu/~bestvina/eprints/questions-updated.pdf}. Reported access: 2026-07-29; extraction: manual transcription. The background says the author PDF presents an open item and current status was not independently verified. Status and rights require review; the PDF is described as publicly author-hosted.
\item The supplied research summary names Brinkmann, Button, Hagen--Wise, and Agol as sources of ingredients, without full bibliographic entries. It expressly says no publication of the explicit Q 12.16 corollary was found and that its author checked the combination of ingredients.
\item Williams' eigenvalue-persistence note, \url{https://arxiv.org/abs/1206.4926}, is identified in the supplied summary as an arXiv preprint and explicitly not used in the reported solution.
\end{itemize}
\end{document}
